\documentclass[11pt]{article}
\usepackage[margin=0.9in]{geometry}
\usepackage{hyperref}
\hypersetup{pdftitle={General Summary: Coordinate-selective resetting},pdfauthor={G. Blake Pierpoint; Olivier Bernard; Yichen Liu},pdfsubject={Plain-language summary of reset observability},pdfkeywords={stochastic resetting, active matter, observability}}
\begin{document}
\section*{Introduction}
Random resetting returns selected parts of a moving particle's state to chosen targets. Returning position, stopping velocity, and realigning direction can affect confinement and motion differently. We compare these effects in one model.

\section*{Main Results}
For an inertial self-propelled particle, we analyze all seven ways to reset position, velocity, and orientation. Position reset localizes without changing internal dynamics. When velocity is retained, orientation reset changes directional order and can suppress or enhance kinetic content depending on rate and inertia. Velocity reset gives one kinetic trace across four protocols. Along a one-parameter curve, velocity-only and orientation-only resets spread identically but drift differently. With known, matched parameters and exact long-time observables, these features identify the reset rule; reduced measurements conceal specific differences.

\section*{Broader Implications}
The work gives a precise example of how interventions leave long-term statistical fingerprints. It may help organize inverse questions in active matter and other noisy systems with one-way coupled state variables. It does not establish finite-data inference or identification with unknown parameters. Within the stated matched-parameter setting, however, the exact decoder shows that confinement, directional order, and kinetic content together retain the full reset-protocol identity.
\end{document}
